% Part: first-order-logic
% Chapter: introduction
% Section: satisfaction

\documentclass[../../../include/open-logic-section]{subfiles}

\begin{document}

\olfileid{fol}{int}{sat}

\olsection{Wanneer een formule waar is (vervulling)}

Nu we weten wat !!{formula}s zijn, kunnen we alvast naar de semantiek
van de eerste-ordelogica kijken. Daarvoor moeten we eerst !!a{structure}
definiëren. In onze eenvoudige taal heeft !!a{structure}~$\Struct M$
drie delen. Het eerste is een niet-lege verzameling $\Domain{M}$, het
\emph{!!{domain}}. Het tweede is het object dat $\Obj a$ in~$\Struct M$
aanwijst. Het derde is de verzameling objecten waarvoor $\Obj P$ in
$\Struct M$ waar is. Het door~$\Obj a$ aangewezen object schrijven we
als $\Assign{\Obj a}{M}$ en de verzameling waarvoor $\Obj P$ waar is
als $\Assign{\Obj P}{M}$. !!^a{structure}~$\Struct{M}$ bestaat dus uit
precies $\Domain{M}$, $\Assign{\Obj a}{M} \in \Domain{M}$ en
$\Assign{\Obj P}{M} \subseteq \Domain{M}$. In de algemene situatie
zijn er veel !!{predicate}s en !!{constant}s, kunnen de !!{constant}s
meer dan één plaats hebben en zijn er ook !!{function}s. Daarom wordt
die situatie ingewikkelder.

Hiermee kunnen we al bepalen wanneer !!{formula}s zonder
!!{variable}s waar zijn. We gebruiken een inductieve definitie die op
de definitie van !!{formula}s lijkt. Eerst leggen we vast wanneer een
atomaire formule waar is in~$\Struct{M}$. Daarna leggen we bijvoorbeeld
vast wanneer $\lnot !A$ waar is in~$\Struct{M}$ door te kijken of $!A$
waar is in~$\Struct{M}$. Dat kan zo:
\begin{enumerate}
  \item $\Atom{\Obj P}{\Obj a}$ is waar in~$\Struct{M}$ precies als
  $\Assign{\Obj a}{M} \in \Assign{\Obj P}{M}$.
  \item $\lnot !A$ is waar in~$\Struct{M}$ precies als $!A$ niet waar
  is in~$\Struct{M}$.
  \item $(!A \land !B)$ is waar in~$\Struct{M}$ precies als $!A$ waar
  is in~$\Struct{M}$ en $!B$ ook waar is in~$\Struct{M}$.
\end{enumerate}
Neem $\Domain{M} = \{0, 1, 2\}$, $\Assign{\Obj a}{M} = 1$ en
$\Assign{\Obj P}{M} = \{1, 2\}$. Dan is
$\Atom{\Obj P}{\Obj a}$ waar in~$\Struct{M}$, want
$\Assign{\Obj a}{M} =  1 \in \{1,2\} = \Assign{\Obj P}{M}$. Daarom is
$\lnot \Atom{\Obj P}{\Obj a}$ niet waar in~$\Struct{M}$. De formule
$\lnot\lnot \Atom{\Obj P}{\Obj a}$ is er wel waar en
$(\lnot \Atom{\Obj P}{\Obj a} \land \Atom{\Obj P}{\Obj a})$ niet.
Zo kunnen we doorgaan.

Bij de kwantoren ontstaat een probleem. We zouden willen zeggen:
``$\lexists[\Obj v_0][\Atom{\Obj P}{\Obj v_0}]$ is waar precies als
$\Atom{\Obj P}{\Obj v_0}$ waar is.'' Maar de
!!{structure}~$\Struct{M}$ zegt niet welke waarde de !!{variable}s
hebben. We willen eigenlijk zeggen dat $\Atom{\Obj P}{\Obj v_0}$ waar
is \emph{voor een of andere waarde van~$\Obj v_0$}. Daarom moeten we
!!{element}s van~$\Domain{M}$ niet alleen aan $\Obj a$, maar ook
aan~$\Obj v_0$ kunnen toekennen. We voeren daarvoor \emph{toekenningen}
aan !!{variable}s in. !!^a{variable} toekenning is een functie~$s$ die
iedere !!{variable} koppelt aan !!a{element} van~$\Domain{M}$, in ons
voorbeeld aan een van $1$, $2$ of~$3$. We weten vooraf niet welke
!!{variable}s in !!a{formula} zullen staan. Daarom kunnen we niet
beperken aan welke !!{variable}s $s$ een waarde geeft. De eenvoudige
oplossing is dat $s$ aan \emph{alle} !!{variable}s~$\Obj v_0$,
$\Obj v_1$, \dots\@ een waarde toekent. We gebruiken later alleen de
waarden die nodig zijn.

Of !!{formula}s waar zijn, hangt nu niet alleen af van
!!a{structure}. Het hangt af van !!a{structure}~$\Struct{M}$ \emph{en}
!!a{variable} toekenning~$s$. We schrijven dit kort als
$\Sat{M}{!A}[s]$. Voor atomaire !!{formula}s met !!{variable}s komt er
een extra regel bij:
\begin{enumerate}
  \item $\Sat{M}{\Atom{\Obj P}{\Obj a}}[s]$ precies als
  $\Assign{\Obj a}{M} \in \Assign{\Obj P}{M}$.
  \item $\Sat{M}{\Atom{\Obj P}{\Obj v_i}}[s]$ precies als
  $s(\Obj v_i) \in \Assign{\Obj P}{M}$.
  \item $\Sat{M}{\lnot !A}[s]$ precies als niet $\Sat{M}{!A}[s]$.
  \item $\Sat{M}{(!A \land !B)}[s]$ precies als $\Sat{M}{!A}[s]$ en $\Sat{M}{!B}[s]$.
\end{enumerate}
Hiermee kunnen we zeggen wanneer $\Struct{M}$ de formule
$\Atom{\Obj P}{\Obj v_0}$ waar maakt bij de waarde~$s(\Obj v_0)$. Voor
$\lexists[\Obj v_0][\Atom{\Obj P}{\Obj v_0}]$ is nog één stap nodig.
We willen dat $\Sat{M}{\lexists[\Obj v_0][\Atom{\Obj P}{\Obj
v_0}]}[s]$ precies als $\Sat{M}{\Atom{\Obj P}{\Obj v_0}}[s']$ voor
\emph{een of andere} manier~$s'$ om aan~$\Obj v_0$ een waarde toe te
kennen. De aan~$\Obj v_0$ toegekende waarde hoeft niet de waarde te zijn
die $s(\Obj v_0)$ aanwijst. We gebruiken daarom een nieuwe notatie. Als
$m \in \Domain{M}$, dan is $\Subst{s}{m}{\Obj v_0}$ bijna dezelfde
toekenning als~$s$. Voor alle !!{variable}s behalve~$\Obj v_0$ zijn ze
gelijk, maar aan $\Obj v_0$ kent de nieuwe toekenning~$m$ toe. De regel
wordt dan:
\begin{enumerate}\setcounter{enumi}{4}
  \item $\Sat{M}{\lexists[\Obj v_i][!A]}[s]$ precies als
  $\Sat{M}{!A}[\Subst{s}{m}{\Obj v_i}]$ voor een~$m \in \Domain{M}$.
\end{enumerate}
Werkt dit? Neem $s(\Obj v_i) = 0$ voor alle~$i \in \Nat$. Dan geldt
$\Sat{M}{\lexists[\Obj v_0][\Atom{\Obj P}{\Obj v_0}]}[s]$ precies als
er een $m \in \Domain{M}$ is waarvoor
$\Sat{M}{\Atom{\Obj P}{\Obj v_0}}[\Subst{s}{m}{\Obj v_0}]$. Zo'n waarde
is er: we kunnen $m = 1$ of $m = 2$ kiezen. De waarde~$s(\Obj v_0)$
die aan~$\Obj v_0$ wordt toegekend door~$s$ zelf is hier $0$, en die
werkt niet. Toch is de formule
waar, want
$\Sat{M}{\Atom{\Obj P}{\Obj v_0}}[\Subst{s}{1}{\Obj v_0}]$ geldt,
maar $\Sat{M}{\Atom{\Obj P}{\Obj v_0}}[s]$ niet.

Dit ziet er verwarrend en omslachtig uit omdat het dat ook is. Toch is
de extra uitwerking nodig. Alleen zo krijgen we een precieze,
inductieve definitie van wanneer alle !!{formula}s waar zijn. Zo'n
definitie hebben we weer nodig om de semantische begrippen precies
vast te leggen. Andere werkwijzen bestaan, maar die zijn allemaal even
(on)elegant.

\end{document}
